% Bagian: computability
% Bab: computability-theory
% Subbagian: russells-paradox

\documentclass[../../../include/open-logic-section]{subfiles}

\begin{document}

\olfileid[id]{cmp}{thy}{rus}
\olsection{Perbandingan dengan Paradoks Russell}

Ada gunanya membandingkan dan membedakan argumen-argumen dalam bagian ini
dengan paradoks Russell:
\begin{enumerate}
\item Paradoks Russell: misalkan $S = \Setabs{x}{x \notin x}$. Maka $S \in S$
  jika dan hanya jika $S \notin S$, suatu kontradiksi.

  \emph{Kesimpulan:} Tidak ada himpunan~$S$ semacam itu. Mengasumsikan
  keberadaan suatu ``himpunan semua himpunan'' tidak konsisten dengan
  aksioma-aksioma teori himpunan lainnya.

\item Modifikasi paradoks Russell: misalkan $F$ merupakan ``fungsi'' dari
  himpunan semua fungsi ke $\{ 0, 1 \}$ yang didefinisikan oleh
  \[
  F(f) =
  \begin{cases}
    1 & \text{jika $f$ berada dalam domain $f$, dan $f(f) = 0$} \\
    0 & \text{selain itu}
  \end{cases}
  \]
  Argumen serupa menunjukkan bahwa $F(F) = 0$ jika dan hanya jika $F(F) = 1$,
  suatu kontradiksi.

  \emph{Kesimpulan:} $F$ bukan suatu fungsi. ``Himpunan semua fungsi'' terlalu
  besar untuk menjadi domain suatu fungsi.

\item Argumen diagonalisasi: misalkan $f_0$, $f_1$, \dots merupakan enumerasi
  fungsi dapat dihitung parsial, dan misalkan $G \colon \Nat \to \{ 0, 1 \}$
  didefinisikan oleh
  \[
  G(x) =
  \begin{cases}
    1 & \text{jika $f_x(x)\downarrow = 0$} \\
    0 & \text{selain itu}
  \end{cases}
  \]
  Jika $G$ dapat dihitung, maka $G$ merupakan fungsi $f_k$ untuk suatu $k$.
  Namun, dalam hal itu $G(k) = 1$ jika dan hanya jika $G(k) = 0$, suatu
  kontradiksi.

  \emph{Kesimpulan:} $G$ tidak dapat dihitung. Perhatikan bahwa berdasarkan
  aksioma teori himpunan, $G$ tetap merupakan suatu fungsi; tidak ada paradoks
  di sini, hanya suatu penjernihan.
\end{enumerate}

Pembicaraan tentang fungsi parsial, fungsi dapat dihitung, fungsi dapat
dihitung parsial, dan seterusnya dapat membingungkan. Himpunan semua fungsi
parsial dari $\Nat$ ke $\Nat$ merupakan koleksi objek yang besar. Sebagian di
antaranya total, sebagian dapat dihitung, sebagian bersifat total sekaligus
dapat dihitung, dan sebagian bukan keduanya. Ingatlah bahwa ketika kita
mengatakan ``fungsi,'' secara bawaan yang kita maksud adalah fungsi total. Jadi,
kita mempunyai:
\begin{enumerate}
\item fungsi dapat dihitung
\item fungsi dapat dihitung parsial yang tidak total
\item fungsi yang tidak dapat dihitung
\item fungsi parsial yang tidak total maupun dapat dihitung
\end{enumerate}
Untuk memilahnya, mungkin berguna untuk menggambar suatu persegi besar yang
mewakili semua fungsi parsial dari $\Nat$ ke $\Nat$, lalu menandai dua daerah
yang saling tumpang tindih, yang masing-masing bersesuaian dengan fungsi total
dan fungsi parsial dapat dihitung. Latihan yang baik ialah mencoba
mendeskripsikan suatu objek di setiap daerah yang dihasilkan dalam diagram itu.

\end{document}

